\begin{Answer}{5.3}
 (a)$x_0=1+(-2)^0 = 1 + 1 = 2$
 (b)$x_1=1+(-2)^1 = 1-2=-1$
 (c)$x_2=1+(-2)^2=1+4= 5$
 (d)$x_3=1+(-2)^3=1-8=-7$
 (e)$x_4=1+(-2)^4=1+16=17$
\end{Answer}
\begin{Answer}{5.4}
We will just provide the final answer for these.  If you can't get these answers, you may need to brush
up on your algebra skills.
(a) $2$, $1/2$,$5/4$, $7/8$, $17/16$;
(b) $2$, $2$, $3$, $7$, $25$;
(c) $1/3$, $1/5$, $1/25$, $1/119$, $1/721$;
(d) $2$, $9/4$, $64/27$, $625/256$, $7776/3125$
\end{Answer}
\begin{Answer}{5.8}
Notice that $x_0=1$, $x_1=5\cdot 1 = 5$, $x_2=5\cdot 5=5^2$, $x_3=5\cdot 5^2 = 5^3$, etc.
Looking back, we can see that $1=5^0$, so $x_0=5^0$.  Also, $x_1=5=5^1$.  So it seems likely that the
solution is $x_n=5^n$.  This is {\em not} a proof, though!
\end{Answer}
\begin{Answer}{5.9}
Notice that $x_0=1$,
$x_1=1\cdot 1 = 1$,
$x_2=2\cdot 1=2$,
$x_3=3\cdot 2 = 6$,
$x_4=4\cdot 6 = 24$,
$x_3=5\cdot 24 = 120$, etc.
Written this way, no obvious pattern is emerging.  Sometimes {\em how} you write the numbers matters.
Let's try this again:
$x_1=1\cdot 1 = 1!$,
$x_2=2\cdot 1=2!$,
$x_3=3\cdot 2\cdot 1=3!$,
$x_4=4\cdot 3\cdot 2\cdot 1=4!$,
$x_3=5\cdot 4\cdot 3\cdot 2\cdot 1=5!$, etc.
Now we can see that $x_n=n!$ is a likely solution.  Again, this isn't a proof.
\end{Answer}
\begin{Answer}{5.10}
Their calculations are correct (Did you check them with a calculator?  You {\em should} have!
How else can you tell whether or not their solution is correct?).
So it does seem like $a_n=2^n$ is the correct solution.
However,
$$\begin{array}{lllllll}
a_5  = \left\lfloor \frac{1+\sqrt{5}}{2}\times a_4\right\rfloor+a_3=
 \left\lfloor \frac{1+\sqrt{5}}{2}\times 16\right\rfloor+8=33\neq 2^5
\end{array}$$
so the solution that seems `obvious' turns out to incorrect.
We won't give the actual solution since the point of this example is to demonstrate
that just because a pattern holds for the first several terms of a sequence, it does not
guarantee that it holds for the whole sequence.
\end{Answer}
\begin{Answer}{5.12}
Hopefully you came up with the solution $x_n=5^n$.
Since $x_0=1=5^0$, it works for the initial condition.
If we plug this back into the
right hand side of $x_n=5\cdot x_{n-1}$,  we get
\begin{eqnarray*}
5\cdot x_{n-1}&=&5\cdot 5^{n-1}\\
&=&5^n\\
&=&x_n,
\end{eqnarray*}
which verifies the formula.  Therefore $x_n=5^n$ is the solution.
\end{Answer}
\begin{Answer}{5.13}
Hopefully you came up with the solution $x_n=n!$.
Since $x_0=1=0!$, it works for the initial condition.
If we plug this back into the
right hand side of $x_n=n\cdot x_{n-1}$,  we get
\begin{eqnarray*}
n\cdot x_{n-1}&=&n\cdot (n-1)!\\
&=&n!\\
&=&x_n,
\end{eqnarray*}
which verifies the formula.  Therefore $x_n=n!$ is the solution.
\end{Answer}
\begin{Answer}{5.14}
The computations are correct, the conclusion is correct, but unfortunately,
the final code has a serious problem.  It works {\em most} of the time,
but it does not deal with negative values correctly.
It should return $3$ for all negative
values, but it continues to use the formula.
The problem is they forgot to even consider what the function does for negative values of $n$.
They probably could have formatted their answer better, too.  It's difficult to follow in paragraph
form.  They could have put the various values of \verb+ferzle(n)+ each on their own line
and presented it mathematically instead of in sentence form.
For instance, instead of `ferzle(1) returns ferzle(0)+2, which is $3+2=5$,'
they should have `$ferzle(1)=ferzle(0)+2=3+2=5$.'
It would have made it much easier to see the pattern.
\end{Answer}
\begin{Answer}{5.15}
\begin{code}
	int ferzle(int n) {
		if(n<=0) {
		    return 3;
	    } else {
		    return 2*n+3;
		}
	}
\end{code}
\end{Answer}
\begin{Answer}{5.19}
We didn't do anything wrong.  We wrote the inequality in the other order, and the indexes are one lower
than those given in the definition.  But that's O.K.  The definition is simply trying to convey the idea
that every term is strictly greater than the previous.  That is what we showed.
We can show that
$x_n<x_{n+1}$, $x_{n+1}>x_n$, $x_{n-1}<x_{n}$, or $x_n>x_{n-1}$.  They all mean essentially the same thing.
The only difference is the order in which the inequalities are written (the first two and last two are
saying exactly the same thing--we just flipped the inequality) and what values of $n$ are valid.
For instance, if the sequence starts at $0$, then we need to assume
$n\geq 0$ for the first pair of inequalities and  $n\geq 1$ for the second pair.
\end{Answer}
\begin{Answer}{5.21}
If you got stuck on this one, first realize that $x_n=\dfrac{n^2 + 1}{n} = n + \dfrac{1}{n}$.
This form might make the algebra a little easier.
Then, follow the technique of the previous example--show that
$x_{n + 1} - x_n >0$.  So, if necessary, go back and try again.
If you already attempted a proof, you may proceed to read the solution.

Notice that, $$ \begin{array}{lll}x_{n + 1} - x_n & = & \left(n + 1 +
\dfrac{1}{n + 1}\right) -\left(n  + \dfrac{1}{n
}\right) \\
& = & 1 + \dfrac{1}{n + 1} - \dfrac{1}{n} \\
& = & 1 -\dfrac{1}{n(n +1)} \\
& > & 0,\\
\end{array}$$
the last step since  $1/n(n+1)<1$ when $n\geq 1$.
Therefore, $x_{n + 1} - x_n > 0$, so $x_{n + 1}> x_n,$
i.e., the sequence is strictly increasing.
If your solution is significantly different than this, make sure you determine one way or another
if it is correct.
\end{Answer}
\begin{Answer}{5.22}
We could go into much more detail than we do here, and hopefully you did when you wrote down
your solutions.  But we'll settle for short, informal arguments this time.
 (a) This is just a linear function.  It is {\bf strictly increasing}.
 (b) Since this keeps going from positive to negative to positive, etc. it is {\bf non-monotonic}.
 (c) We know that $n!$ is strictly increasing.  Since this is the reciprocal of that function, it
 is almost {\em strictly decreasing} (since we are dividing by a number that is getting larger).
 However, since $1/0!=1/1!=1$, it is just {\bf decreasing}.
(d) This is getting closer to $1$ as $n$ increases.  It is {\bf strictly increasing}
(e) This is $n(n-1)$.  $x_1=0$, $x_2=2$, $x_3=6$, etc.
Each term is multiplying two numbers that are both
getting larger, so it is {\bf strictly increasing}.
(f) This is similar to the previous one, but $x_0=x_1=0$, so it is just {\bf increasing}.
(g) This alternates between $-1$ and $1$, so it is {\bf non-monotonic}.
(h) Each term subtracts from $1$ a smaller numbers than the last term, so it is {\bf strictly increasing}.
(i) Each term adds to $1$ a smaller number than the last term, so it is {\bf strictly decreasing}.
\end{Answer}
\begin{Answer}{5.26}
You should have concluded that $a=-\frac{2}{3^{17}}$ and that
$r=\frac{2}{3^{16}}/\left(-\frac{2}{3^{17}}\right)=-3^{17}/3^{16}=-3$
(or you could have divided the second and third terms).
Then the $n$-th term is
$-\frac{2}{3^{17}}\left(-3\right)^{n-1}=\frac{2(-1)^{n}}{3^{18-n}}$
(Make sure you can do the algebra to get to this simplified form).
Finally, the $17$th term is
$\frac{2(-1)^{17}}{3^{18-17}}=-\dfrac{2}{3}$
\end{Answer}
\begin{Answer}{5.28}
We are given that $ar^5 = 20$ and $ar^9 = 320$.
Dividing, we can see that $r^4=16$.  Thus $r=\pm 2$.
(We don't have enough information to know which it is).
Since $ar^5=20$, we know that $a=20/r^5=\pm 20/32$.
So the third term is $ar^2=(\pm 20/32)(\pm 2)^2=\pm 80/32 = \pm 5/2$.
Thus $|ar^2|=5/2$.
\end{Answer}
\begin{Answer}{5.32}
\begin{lenum}
\item
The difference between the each of the first 4 terms of the sequence is 7, so it
appears to be an arithmetic sequence.  Doing a little math, the correct answer
appears to be (d) 51.
\item
Although the sequence {\em appears to be} arithmetic, we cannot be certain that it is.
If you are told it is arithmetic, then 51 is absolutely the correct answer.
Notice that the previous example specifically stated that you should assume that the
pattern continues.  This one did not.
Without being told this, the rest of the sequence could be anything.
The $8$th term could be $0$ or $8,675,309$ for all we know.
Of the choices given, 51 is the most {\em obvious} choice, but any of the answers
could be correct.  This is one reason I hate these sorts of questions on tests.

Although I think it is important to point out the flaw in these sorts of questions,
it is also important to conform to the expectations when answering such questions
on standardized tests.
In other words, instead of disputing the question (as some students might be inclined
to do), just go with the obvious interpretation.
\end{lenum}
\end{Answer}
\begin{Answer}{5.33}
(a) The closed form was $x_n=5^n$, which is clearly geometric (with $a=1$ and $r=5$) and not arithmetic.
(b) Since the solution for this one is $x_n=n!$, this is neither arithmetic or geometric.
(c) Since the sequence is essentially $f_n=2n+3$, with initial condition $f_0=3$, it is an
arithmetic sequence.  It is clearly not geometric.
\end{Answer}
\begin{Answer}{5.36}
$\dsum_{i=0}^{100} y^i$
\end{Answer}
\begin{Answer}{5.38}
$\dsum_{i=0}^{50} (y^2)^i \ \ \ \ \ \mbox{ or } \ \ \ \ \ \dsum_{i=0}^{50} y^{2i}$
\end{Answer}
\begin{Answer}{5.40}
(a) $2$
(b)$11$
(c) $100$
(d) $101$
\end{Answer}
\begin{Answer}{5.45}
(a) ${\displaystyle \sum_{k=5}^{6} 5 = 5 \sum_{k=5}^{6} 1= 5\cdot 2 = 10 }$.
(b) ${\displaystyle \sum_{k=20}^{30} 200 = 200\sum_{k=20}^{30} 1=200(30-20+1)=2200}$.
\end{Answer}
\begin{Answer}{5.48}
Using Theorem~\ref{thm:sumconstant}, we get the following answers:
%(a) $2\cdot 5 = 10$.
(a) $(30-20+1)200=11*200=2200$.
(b) $900$
(c) $909$. Notice that this one has one more term than the previous one.
The fact that the additional index is $0$ doesn't matter since it is adding $9$ for that term.
\end{Answer}
\begin{Answer}{5.49}
This solution contains an `off by one' error.  The correct answer is $10(75-25+1)=10*51=510$.
\end{Answer}
\begin{Answer}{5.52}
(a) $20\cdot 21/2=210$
(b) $100\cdot 101/2=5050$
(c) $1000\cdot 1001/2=500500$
\end{Answer}
\begin{Answer}{5.53}
\begin{solevalenum}
\item
Another example of the `off by one error'.  They are using the formula $n(n-1)/2$
instead of $n(n+1)/2$.
\item
This answer doesn't even make sense.  What is $k$ in the answer?  $k$ is just an
index of the summation.  The index should {\em never} appear in the answer.
The problem is that you can't pull the $k$ out of the sum since each term in the
sum depends on it.
\end{solevalenum}
\end{Answer}
\begin{Answer}{5.54}
It is true.  The additional term that the sum adds is $0$, so the sum is the same whether
or not it starts at $0$ or $1$.
\end{Answer}
\begin{Answer}{5.57}
$\dsum_{i=1}^{100} 2 - i =\dsum_{i=1}^{100} 2 - \dsum_{i=1}^{100} i
=200 - \dfrac{100\cdot 101}{2} = 200-5050 = -4850.$
\end{Answer}
\begin{Answer}{5.58}
The sum of the first $n$ odd integers is

\noindent
$\dsum_{k=1}^{n}(2k-1)=
\dsum_{k=1}^{n}2k-\dsum_{k=1}^{n}1=
2\dsum_{k=1}^{n}k-\dsum_{k=1}^{n}1=
2\frac{n(n+1)}{2}-n=n^2+n-n=n^2.
$
\end{Answer}
\begin{Answer}{5.61}
(a) ${\displaystyle \sum_{k=10}^{20} k
=\sum_{k=1}^{20} k - \sum_{k=1}^{9} k
=20\cdot 21/2 - 9\cdot 10/2 = 210 - 45 = 165}$.

(b) ${\displaystyle \sum_{k=21}^{40} k
=\sum_{k=1}^{40} k - \sum_{k=1}^{20} k
=40\cdot 41/2 - 20\cdot 21/2 = 820-210 = 610}.$
\end{Answer}
\begin{Answer}{5.62}
\begin{solevalenum}
\item
Another example of the off-by-one error.
The second sum should end at $29$, not $30$.
\item
This one has two errors, one of which is repeated twice.
It has the same error as the previous solution, but it also uses
the incorrect formula for each of the sums (the off-by-one error).
\item
This one is correct.
\end{solevalenum}
\end{Answer}
\begin{Answer}{5.63}
Two errors are made that cancel each other out.
The first error is that the second sum in the second step should go to $29$, not $30$.
But in the computation of that sum in the next step,
the formula $n(n-1)/2$ is used instead of $n(n+1)/2$
(The correct formula was used for the first sum).
This is a rare case where an off-by-one error is followed by the opposite off-by-one error
that results in the correct answer.

It should be emphasized that even though the correct answer is obtained,
{\em this is an incorrect solution}.
They obtained the correct answer by sheer luck.
\end{Answer}
\begin{Answer}{5.65}
There are two ways to answer this.  The smart aleck answer is `because it is correct.'
But {\em why} is it correct with $2$, and couldn't it be slightly modified to work with $1$ or $0$?
The answer is {\em no} because if you plug $1$ or $0$ into $\frac{1}{(k-1)k}$,
you get a division by $0$.  Hopefully I don't need to tell you that this is a bad thing.
\end{Answer}
\begin{Answer}{5.66}
\begin{eqnarray*}
\dsum _{k=1}^n k^3+k &=& \dsum _{k=1}^n k^3+ \dsum _{k=1}^n k\\
&=&\frac{n^2(n+1)^2}{4}+\frac{n(n+1)}{2}\\
&=&\frac{n(n+1)}{2}\left(\frac{n(n+1)}{2}+1\right)\\
&=&\frac{n(n+1)}{2}\left(\frac{n^2+n+2}{2}\right)\\
&=&\frac{n(n+1)(n^2+n+2)}{4}
\end{eqnarray*}
\end{Answer}
\begin{Answer}{5.68}
(a)
$\dsum _{i=1}^n \dsum _{j=1}^i 1 = \dsum _{i=1}^n i = \dfrac{n(n+1)}{2}.$

\noindent
(b) $\dsum _{i=1}^n \ \dsum _{j=1}^i\ j
= \dsum _{i=1}^n \dfrac{i(i+1)}{2} = \cdots = \dfrac{n(n+1)(n+2)}{6}.$
(This one involves doing a little algebra, applying two formulas, and then doing a little more
algebra.  Make sure you work it out until you get this answer.)

\noindent
(c)
$\dsum _{i=1}^n \ \dsum _{j=1}^n\ ij=
\dsum _{i=1}^n\left( i\dsum _{j=1}^n j\right)=
\dsum _{i=1}^n\left( i\dfrac{n(n+1)}{2}\right)=
\dfrac{n(n+1)}{2}\dsum _{i=1}^n i=
\dfrac{n(n+1)}{2}\dfrac{n(n+1)}{2}=
\dfrac{n^2(n+1)^2}{4}.$
\end{Answer}
\begin{Answer}{5.72}
$\frac{3^{50} - 1}{2}=358948993845926294385124$.
\end{Answer}
\begin{Answer}{5.73}
This is equivalent to $\sum_{k=0}^{34}(-2)^k$,
 so the summation is
 $(1-(-2)^{35})/(1-(-2))=(1-(-1)^{35}2^{35})/3=(1+2^{35})/3=11453246123$.
\end{Answer}
\begin{Answer}{5.74}
 (a) $\frac{1 - y^{101}}{1 - y}$ or $\frac{y^{101}-1}{y-1}$ (We won't give the alternatives for the rest.
 If your answer differs, do some algebra to make sure it is equivalent.)
 (b) $\frac{1 -(-y)^{101}}{1 - (-y)} = \frac{1 + y^{101}}{1 + y}$
 (c) $\frac{1 - y^{102}}{1 - y^2}$.
\end{Answer}
\begin{Answer}{5.77}
$x^5-1=(x-1)(x^4+x^3+x^2+x+1)$.
\end{Answer}
\begin{Answer}{5.78}
$2^1+2^2 + 2^3  + \cdots + 2^{n+1}$;
$2^0$;
$2^{n+1}$;
$2^{n+1}-2^0$
\end{Answer}
\begin{Answer}{5.80}
$a\dsum_{k=0}^{n}r^k$;
$a\frac{1-r^{k+1}}{1-r}$.
\end{Answer}
\begin{Answer}{5.81}
Let $S = a + ar + ar^2 + \cdots + ar^{n}.$
Then $rS = ar + ar^2 + \cdots + ar^{n+1}Z$, so
\begin{eqnarray*}
S - rS &=& a + ar + ar^2 + \cdots + ar^{n} - ar - ar^2 - \cdots - ar^{n+1}\\
&=& a - ar^{n+1}.
\end{eqnarray*}
From this we deduce that $$S = \frac{a - ar^{n+1}}{1 - r},$$
that is,
$$ \sum_{k=0}^{n}ar^k= \frac{a - ar^{n+1}}{1 - r}.$$
\end{Answer}
\begin{Answer}{5.92}
 $A = \begin{bmatrix} 1 & 1 & 1 \cr
2 & 4 & 8 \cr 3 & 9 & 27 \cr \end{bmatrix}$
\end{Answer}
\begin{Answer}{5.96}
${\bf 0}_{3\times 4}=
\begin{bmatrix}0 & 0 & 0 & 0 \cr
0 & 0 & 0 & 0 \cr
0 & 0 & 0 & 0 \cr \end{bmatrix}$
and
 ${\bf I}_5=
\begin{bmatrix}
1 & 0 & 0 & 0 & 0\cr
0 & 1 & 0 & 0 & 0\cr
0 & 0 & 1 & 0 & 0\cr
0 & 0 & 0 & 1 & 0\cr
0 & 0 & 0 & 0 & 1\cr\end{bmatrix}$
\end{Answer}
\begin{Answer}{5.99}
$\displaystyle A + 2B = \begin{bmatrix} -1 & 3 \cr 3 & 3 \cr 0 & 0
\end{bmatrix}$
\end{Answer}
\begin{Answer}{5.100}
$M + N = \begin{bmatrix}a + 1 & 0 & 2c \cr  a &  b - 2a & 0 \cr 2a
& 0 & -2 \cr\end{bmatrix}, \ \ \ 2M = \begin{bmatrix}2a & -4a & 2c
\cr 0 & -2a & 2b \cr 2a + 2b & 0 & -2 \cr\end{bmatrix}$
\end{Answer}
\begin{Answer}{5.102}
 $x = 1$ and $y = 4$.  
\end{Answer}
\begin{Answer}{5.105}
 $\begin{array}{llll}AA & = &
\begin{bmatrix} 2 & 1 & 3 \cr
 0 & 1 & 1 \cr
  4 & 4 & 0 \cr
  \end{bmatrix}\begin{bmatrix} 2 & 1 & 3 \cr
 0 & 1 & 1 \cr
  4 & 4 & 0 \cr
  \end{bmatrix} \vspace{2mm}
   =  \begin{bmatrix} 16 & 15 & 7 \cr
 4 & 5 & 1 \cr
  8 & 8 & 16 \cr
  \end{bmatrix}
 \end{array} $
\end{Answer}
\begin{Answer}{5.107}
$AB = \begin{bmatrix} a & b & c \cr c+a & a + b & b + c
\cr a + b + c & a + b + c & a + b + c \cr
 \end{bmatrix}, \ BA = \begin{bmatrix} a + b + c & b + c & c
 \cr a + b + c & a + b & b \cr a + b + c & c + a & a \cr
\end{bmatrix}
 $
 
\end{Answer}
\begin{Answer}{5.110}
 $ \begin{bmatrix} 32 & -32  \cr -32 & 32 \cr\end{bmatrix}$.
\end{Answer}
\begin{Answer}{5.115}
Observe that
$A^2 = (AB)(AB) = A(BA)B = A(B)B = (AB)B = AB =A. $
Similarly,
$B^2 = (BA)(BA) = B(AB)A = B(A)A = (BA)A = BA =B. $
\end{Answer}
\begin{Answer}{5.116}
Disprove! Take $\dis{A = \left[\begin{array}{ll}1 & 0 \\ 0 & 0
\end{array}\right]}$ and $\dis{B = \left[\begin{array}{ll}0 & 1 \\ 0 & 0
\end{array}\right]}$. Then $AB = B$, but
$BA = {\bf 0}_2$.
\end{Answer}
\begin{Answer}{5.120}
 We have $\tr{A^2} =
\tr{\begin{bmatrix} a & b \cr c & d \cr\end{bmatrix}\begin{bmatrix}
a & b \cr c & d \cr\end{bmatrix}} = \tr{\begin{bmatrix} a^2 + bc &
ab + bd \cr ca + cd & d^2 + cb \cr\end{bmatrix}} = a^2 + d^2 + 2bc
$
and $\left(\tr{\begin{bmatrix} a
& b \cr c & d \cr\end{bmatrix}}\right)^2 = (a+d)^2.  $ Thus $
\tr{A^2} = (\tr{A})^2 \iff a^2 + d^2 + 2bc = (a + d)^2 \iff bc =
ad,
$
is the condition sought.
\end{Answer}
\begin{Answer}{5.123}
$\displaystyle N^T = \begin{bmatrix}
1&5&9\cr
2&6&10\cr
3&7&11\cr
4&8&12\cr\end{bmatrix}.$
\end{Answer}
\begin{Answer}{5.127}
 We have
$(AB - BA)^T = (AB)^T - (BA)^T = B^TA^T - A^TB^T = -BA - A(-B) = AB - BA.$
\end{Answer}
